\documentclass[11pt, a4paper]{article}
\usepackage{geometry}
\usepackage[parfill]{parskip}
\usepackage{graphicx}
\usepackage{amsmath}
\usepackage{enumerate}

\usepackage{charter}

\title{Huiswerk 2 \\ IB0402 Logica, verzamelingen en relaties}
\author{Harman Singh}
\date{Month Day, Year}

\begin{document}
\maketitle

\pagebreak
\section*{Opgave 4}
\textbf{A}

$M$ = miauwen

$x$ = een katje

\[
\forall x M(x)
\]

vertaling: Voor alle $x$, $x$ miauwt


\textbf{B}

$S$ = scheefhangen 

$x$ = een gordijn

\[
\exists x S(x)
\]

vertaling: Er bestaat een $x$ zodat $x$ scheefhangt



\pagebreak
\section*{Opgave 5}

\textbf{A}

Eerst herschrijf ik de implicatie van form $A \rightarrow B$ naar $\neg A \lor B$:

\[
\forall x P(x) \rightarrow Q(y) \equiv \neg ( \forall x P(x)) \lor Q(y)
\]

Hierna herschrijf ik het linkerdeel met behulp van de kwantorregel (als niet alle $x$ aan $P(x)$ voldoen, dan is er een $x$ die niet aan $P(x)$ voldoet):

\[
\neg ( \forall x P(x)) \equiv \exists x \neg P(x)
\]

En dan heb je de prenex normaalvorm:

\[
\exists x ( \neg P(x) \lor Q(y) )
\]


\textbf{B}

Hier kan ik voor de eerste paar stappen dezelfde regels gebruiken als in deel A.

Dus eerste de implicatie van form $A \rightarrow B$ herschrijven naar $\neg A \lor B$:

\[
\forall x P(x) \rightarrow Q(x) \equiv \neg ( \forall x P(x)) \lor Q(x)
\]

En dan het linkerdeel herschrijven met behulp van de kwantorregel:

\[
\neg ( \forall x P(x)) \equiv \exists x \neg P(x)
\]

Dan krijg je:

\[
\exists x ( \neg P(x) \lor Q(x) )
\]
Maar hier kan ik de variabele $x$ niet opnieuw gebruiken, omdat die al gebonden is aan de kwantor $\exists x$.
Dus ik hernoem ($\alpha$-renaming) de variabele $x$ in $z$ in het linkerdeel om te bekomen aan de prenex normaalvorm:
\[
\exists z ( \neg P(z) \lor Q(x) )
\]



\pagebreak
\section*{Opgave 6}

\ldots


\end{document}